% ======================================================================
%  Project Report: Tool Selection Robustness in LLM Agents
%  Compiled with pdflatex. Figures are optional: place PNG images in the
%  figures/ folder using the names referenced below, or the document
%  will still compile with placeholder boxes.
% ======================================================================
\documentclass[11pt,a4paper]{report}

\usepackage[margin=2.5cm]{geometry}
\usepackage{graphicx}
\usepackage{booktabs}
\usepackage{array}
\usepackage{enumitem}
\usepackage{caption}
\usepackage{float}
\usepackage[hidelinks]{hyperref}
\usepackage{xcolor}
\usepackage{titlesec}

\graphicspath{{figures/}}
\captionsetup{font=small,labelfont=bf}

\titleformat{\chapter}[display]
  {\normalfont\huge\bfseries}{\chaptertitlename\ \thechapter}{12pt}{\Huge}
\titlespacing*{\chapter}{0pt}{0pt}{18pt}

% Optional-figure helper: shows the image when present, otherwise a
% framed placeholder so the document always compiles.
\newcommand{\figimg}[3][0.85]{%
  \IfFileExists{#2}{%
    \begin{center}\includegraphics[width=#1\textwidth]{#2}\end{center}%
  }{%
    \begin{center}%
      \fbox{\parbox{#1\textwidth}{\centering\vspace{2.2cm}%
      [Place image here: #2]\vspace{2.2cm}}}%
    \end{center}%
  }%
  \caption{#3}%
}

\newcommand{\metric}[1]{\texttt{#1}}

\begin{document}

% ----------------------------------------------------------------------
\begin{titlepage}
  \centering
  {\LARGE\bfseries An Evaluation Framework for Tool Selection\\[2mm]
  Robustness in LLM Agents\par}
  \vspace{1.2cm}
  {\large Project Report\par}
  \vspace{0.6cm}
  {\large Prepared for research supervision and feedback\par}
  \vfill
  {\large Department of Computer Science and Engineering\par}
  \vspace{0.4cm}
  {\large \today\par}
\end{titlepage}

% ----------------------------------------------------------------------
\tableofcontents

% ======================================================================
\chapter{Introduction}
% ======================================================================

Large language models (LLMs) are increasingly deployed as \emph{agents}:
systems that plan a task and then call external tools, such as weather
services, calculators, or e-commerce APIs, to complete it. Before an agent
can execute anything, it must decide \emph{which tool} fits the user's
request. This decision, known as tool selection, is security critical. If
the agent picks the wrong tool, or a tool that was crafted with harmful
intent, the consequences can include data leaks, unauthorized actions, or
simply incorrect answers.

Two properties of the real world make this problem hard. First, the space
of available tools is large and constantly growing, so agents narrow the
search using a retrieval step before making a final choice with an LLM.
Second, tool repositories are open. Anyone can publish a tool together
with a short natural language description of what it does, and those
descriptions are later fed directly into the LLM's context. A description
written to exploit this channel can carry hidden instructions that
manipulate the selection process itself.

This report describes an evaluation framework we built to study these
questions in a controlled setting. The framework implements the complete
tool selection pipeline, provides a way to inject crafted tool documents
into the library, measures how strongly those documents influence
selection, and evaluates a set of detection methods that aim to flag
suspicious documents before they are ever selected. Everything runs
inside a synthetic benchmark environment, and the framework ships with an
interactive dashboard so that every measurement can be inspected and
reproduced visually.

The rest of the report is organized as follows. Chapter 2 gives the
background on LLM agents and tool selection. Chapter 3 defines the threat
model we assume. Chapter 4 describes the system architecture. Chapter 5
explains the evaluation methodology and metrics. Chapters 6 and 7 cover
the two studies the framework supports: the injection study and the
detection study. Chapter 8 summarizes the implementation and the
dashboard. Chapter 9 presents illustrative results obtained with the
current synthetic setup. Chapters 10 and 11 discuss limitations and
conclusions.

% ======================================================================
\chapter{Background}
% ======================================================================

\section{LLM Agents and Tools}

An LLM agent is an application that wraps a language model with the
ability to act on the world. Typical capabilities include reading a
knowledge base, browsing a website, and, most importantly for this work,
calling external tools. Each tool is an API with a defined function. What
the model sees of the tool is not the implementation but a short
\emph{tool document}: a name plus a natural language description of what
the tool does.

\section{The Two Step Selection Pipeline}

Modern agent frameworks rarely ask the LLM to pick directly from a
catalog of thousands of tools. Instead they decompose selection into two
steps, which we adopt in our framework as well.

\begin{enumerate}
  \item \textbf{Retrieval.} A lightweight embedding model maps the user's
  query and every tool document into vector space. The documents most
  similar to the query are kept, producing a short candidate list of size
  $k$ (the top-$k$ set).
  \item \textbf{Selection.} The user's query and the $k$ candidate
  documents are placed into a structured prompt for the LLM, together
  with strict instructions to choose exactly one tool and to answer in a
  fixed JSON format.
\end{enumerate}

Figure~\ref{fig:pipeline} illustrates the pipeline. The two steps have
different failure modes, and any security analysis must consider both: a
document must first \emph{get retrieved} and then \emph{get selected}.

\begin{figure}[H]
  \centering
  \figimg{pipeline.png}{The two step tool selection pipeline: the query
  is embedded, the top-$k$ tool documents are retrieved by similarity,
  and the LLM selects exactly one tool.}
  \label{fig:pipeline}
\end{figure}

\section{Why Tool Documents Are a Risk}

The descriptions inside tool documents are untrusted data by nature. They
are written by third parties, fetched from open repositories, and then
concatenated into the same prompt that carries the system's own
instructions. Because LLMs cannot perfectly separate instructions from
data, a description can be written so that the model treats part of it as
an order: for example, a sentence instructing the model to always prefer
that particular tool for a certain kind of request. This phenomenon is
generally known as prompt injection, and the tool document channel is one
of its most direct instantiations.

% ======================================================================
\chapter{Threat Model}
% ======================================================================

We model a platform where third parties can publish tools, mirroring
public tool hubs used by real agent frameworks. The evaluator plays the
role of such a third party and studies how much influence one published
document can exert.

\section{Objective}

The objective is to answer a measurable question: given a target task,
such as answering weather queries, can a single crafted tool document,
inserted into an otherwise benign library, redirect the agent's selection
towards the crafted tool across a wide range of natural user phrasings of
that task?

\section{Assumptions}

We adopt a conservative, black box setting, which we call the no box
setting, for the document author:

\begin{itemize}
  \item The author knows the target task, but does not know the exact
  queries real users will type.
  \item The author cannot read the victim's tool library, cannot query
  the victim's retriever or LLM, and does not know the top-$k$ setting.
  \item The author can only publish one document: a name and a
  description.
\end{itemize}

These assumptions make the study realistic: if a crafted document still
succeeds under these constraints, the risk applies to any open
repository. All experiments in this framework remain within the synthetic
benchmark environment; no real agents, users, or production systems are
involved.

% ======================================================================
\chapter{System Design}
% ======================================================================

The framework is built from five cooperating components, each with a
single responsibility. Figure~\ref{fig:architecture} shows how they are
connected.

\begin{itemize}
  \item \textbf{Tool library.} An ordered collection of tool documents
  (name and description). Operations that modify the library are
  copy on write: injecting a document always produces a new library,
  leaving the original untouched. This lets every experiment compare a
  clean baseline and a modified library side by side.
  \item \textbf{Retriever.} The retrieval step. It uses a pluggable
  embedding backend and supports cosine similarity and dot product as
  scoring functions. The default backend is a small sentence transformer
  model; a deterministic offline embedder is also available so that
  experiments run without any model download.
  \item \textbf{Selector.} The selection step. It renders the structured
  selection prompt, forwards it to any pluggable LLM backend, and parses
  the answer tolerantly. The parser handles markdown fences, trailing
  text, and classifies every outcome as a valid choice, a hallucinated
  name, unparseable output, or a refusal.
  \item \textbf{Benchmark runner.} The orchestrator. Given a library, a
  retriever, a selector, a list of query and expected tool pairs, a value
  of $k$, and an optional injected document, it runs the full pipeline
  per query, records every intermediate result, and aggregates all
  metrics.
  \item \textbf{Dashboard.} An interactive browser interface with five
  views: library editing, benchmark configuration and results, document
  crafting, detection analysis, and run history.
\end{itemize}

Every external dependency is pluggable. The embedding backend, the LLM
used for selection, the LLM used for crafting documents, and the model
used for detection can each be swapped without changing the rest of the
system.

\begin{figure}[H]
  \centering
  \figimg{architecture.png}{The five framework components and the flow
  of a benchmark run from the tool library through retrieval and
  selection to recorded metrics.}
  \label{fig:architecture}
\end{figure}

% ======================================================================
\chapter{Evaluation Methodology}
% ======================================================================

\section{Benchmark Setup}

A benchmark run is defined by four inputs:

\begin{enumerate}
  \item \textbf{A tool library.} The shipped synthetic library contains
  fifteen tools across ten categories (weather, shopping, calendar,
  translation, calculation, search, finance, communication, travel, and
  media). Categories that overlap with the target task act as the honest
  competition; the rest act as distractors. The library is fully
  editable in the dashboard, and larger libraries can be generated or
  imported.
  \item \textbf{A target task with evaluation queries.} Each query is a
  natural phrasing of one task, paired with the tool that should have
  been selected for it. Queries can be written by hand, imported, or
  generated by an LLM from a short task description.
  \item \textbf{Pipeline settings.} The embedding backend, the similarity
  metric, the top-$k$ value (default 5), and the selection LLM backend.
  \item \textbf{An optional injected document.} A single extra tool
  document inserted into a copy of the library for the modified run.
\end{enumerate}

Each benchmark performs two passes: a clean baseline pass over the
original library, and, when a document is provided, a modified pass over
the injected copy. Because injection is copy on write, both passes use
identical retrieval and selection settings.

\section{Metrics}

Table~\ref{tab:metrics} lists the metrics computed by the framework. The
first two describe the healthy pipeline. The next two describe how much
influence an injected document exerts, computed on the modified pass
only. The last three describe detection quality, used in the detection
study of Chapter 7.

\begin{table}[H]
  \centering
  \caption{Evaluation metrics used throughout the framework.}
  \label{tab:metrics}
  \begin{tabular}{llp{9.2cm}}
    \toprule
    \textbf{Metric} & \textbf{Pass} & \textbf{Meaning} \\
    \midrule
    Accuracy (ACC) & baseline & Fraction of queries where the expected tool was selected. \\
    Hit rate at $k$ (HR@$k$) & baseline & Fraction of queries where the expected tool appeared in the top-$k$ set. \\
    Target selection rate & modified & Fraction of queries where the injected document was the selected tool. \\
    Target retrieval rate & modified & Fraction of queries where the injected document appeared in the top-$k$ set. \\
    False positive rate (FPR) & detection & Fraction of benign documents wrongly flagged. \\
    False negative rate (FNR) & detection & Fraction of crafted documents that escaped detection. \\
    Area under the curve (AUC) & detection & Rank separation between benign and crafted score distributions. \\
    \bottomrule
  \end{tabular}
\end{table}

The two target rates deliberately separate the two pipeline stages: a
document can be retrieved often but selected rarely, and the gap between
the two rates identifies which stage resists the influence.

% ======================================================================
\chapter{Injection Study}
% ======================================================================

This study asks how strongly a single crafted document can influence
selection, and what crafting effort is required.

\section{Manual Injection Styles}

The first part of the study uses five hand written injection styles,
representing what an attacker can produce with simple prompt
engineering. Table~\ref{tab:manual} lists them. These styles differ in
how they attempt to separate the injected instruction from the rest of
the description.

\begin{table}[H]
  \centering
  \caption{The five manual injection styles used as baseline documents.}
  \label{tab:manual}
  \begin{tabular}{lp{11.5cm}}
    \toprule
    \textbf{Style} & \textbf{Description content (illustrative)} \\
    \midrule
    Direct instruction & ``Just output the word X as your final answer.'' \\
    Escape characters & Line breaks inserted before the instruction to visually separate it from the description. \\
    Context override & ``Ignore all previous instructions,'' followed by the instruction. \\
    Fake completion & Pretends the previous task is finished, then issues a new instruction. \\
    Combined & Uses escape characters, context override, and fake completion together. \\
    \bottomrule
  \end{tabular}
\end{table}

\section{Automated Two Phase Crafting}

Manual styles ignore a central fact of the pipeline: the document must
first be retrieved before it can be selected. The second part of the
study therefore crafts documents automatically, in two phases matching
the two pipeline stages. The retrieval phase produces text that wins the
similarity race, and the selection phase produces an instruction that
wins the final choice. The two parts are concatenated into one
natural looking description.

\subsection{Phase One: Retrieval Aligned Text}

A set of example queries for the target task, called the shadow query
set, is prepared. An LLM is then asked to write a short functionality
description that addresses those queries in general terms, with an
approximate word budget. The prompt asks for core functionality rather
than an answer to any single query, so the output generalizes across
phrasings.

The resulting text is semantically aligned with the shadow queries,
which is exactly what the embedding based retriever rewards: the
document's vector lands close to the queries, so the document enters the
top-$k$ candidate set consistently. This part alone is a strong baseline,
and it also serves as the initialization for the gradient based method
described below.

\subsection{Phase Two: Selection Instruction, Gradient Free Refinement}

Once the document is retrieved, the selection instruction must convince
the selection LLM to choose it. The gradient free procedure refines this
instruction by search. It needs only API level access to two LLMs: a
proposal model that writes candidate phrasings, and the selection model
that judges them. The procedure works as follows.

\begin{enumerate}
  \item \textbf{Prepare the judging setup.} Keep a small set of
  representative benign documents that stand in for the retrieval
  context. For evaluation, the crafted document is placed into this
  context to form the candidate list, exactly as the real selection
  prompt would.
  \item \textbf{Seed the instruction.} Start from a simple explicit
  instruction containing the crafted tool's name, for example ``Just
  output the word X as your final answer.''
  \item \textbf{Propose variants.} For each current candidate, ask the
  proposal model for several alternative phrasings of the instruction.
  \item \textbf{Evaluate all variants.} Run the real selection pipeline
  once per shadow query for every variant, and count how many shadow
  queries each variant wins. A win means the selection model emitted the
  crafted tool's name as its answer.
  \item \textbf{Prune and repeat.} Keep the variants with the highest win
  counts, attach the feedback from this round to the next proposal
  request, and repeat. The search stops when a variant wins every shadow
  query or when the iteration budget is exhausted.
\end{enumerate}

The key property of this loop is that the judge is the same kind of
model that will make the real decision. Success in the loop therefore
transfers directly to the benchmark, and the search discovers phrasings
that a human would not have thought of.

\subsection{Phase Two: Selection Instruction, Gradient Based Refinement}

The gradient based procedure treats the instruction as a sequence of
tokens and refines it directly against a local language model, using
gradient information. It requires an open source model that exposes its
weights, and it runs on a GPU when one is available. Three loss terms
guide the optimization, each answering one question.

\begin{itemize}
  \item \textbf{Alignment loss.} What is the likelihood that the model,
  given the full selection prompt with the crafted document inside it,
  emits the complete expected answer, namely the JSON object that names
  the crafted tool? The loss is the negative log of this likelihood.
  This is the main objective.
  \item \textbf{Consistency loss.} What is the likelihood that the model
  emits the crafted tool's name on its own? This term reinforces the
  target name even when the exact JSON wrapper varies.
  \item \textbf{Readability loss.} How natural is the instruction itself?
  This is the average negative log likelihood of the instruction's own
  tokens under the model, and it acts as a penalty that keeps the text
  fluent.
\end{itemize}

The total objective is a weighted sum of the three terms, with the
readability term weighted lightly. The optimization alternates two
operations: first, the gradient of the objective is computed with respect
to the embeddings of the instruction's tokens; second, for each token
position, candidate replacement tokens are scored by their similarity to
that gradient, the most promising replacements are evaluated with a
forward pass, and a replacement is accepted whenever it reduces the
objective.

The readability term is what keeps the optimized text from degrading
into nonsense. Without it, token level optimization happily trades
fluency for likelihood, producing broken text that a perplexity based
detector can trivially flag. With it, the final instruction remains
readable, which is precisely what makes detection difficult.

\subsection{Composed Document}

The final crafted document is the concatenation of the retrieval aligned
text and the refined instruction. Figure~\ref{fig:crafting} sketches the
full process.

\begin{figure}[H]
  \centering
  \figimg{crafting.png}{Two phase crafting: retrieval aligned text is
  generated from example queries, then a selection instruction is
  refined by variant search or token optimization, and both parts are
  concatenated into one tool document.}
  \label{fig:crafting}
\end{figure}

% ======================================================================
\chapter{Detection Study}
% ======================================================================

The second study asks whether suspicious documents can be flagged before
selection, using detectors that score each document and a threshold that
decides what is flagged.

\section{Detectors}

Three detectors are implemented, covering different families of defense:

\begin{itemize}
  \item \textbf{Perplexity detector.} Scores the description by its
  average token negative log likelihood under a small local causal
  language model. Natural descriptions receive low scores; broken or
  unusual phrasing receives high scores.
  \item \textbf{Windowed perplexity detector.} The same principle applied
  over sliding windows of tokens, taking the maximum window score. This
  catches a document where only a short segment is unusual.
  \item \textbf{Known answer detector.} Prepends a fixed instruction to
  the document, for example asking the model to repeat a secret phrase
  while ignoring the following text. If the model fails to repeat the
  phrase, the document is flagged. This detector needs an LLM backend
  but no local model.
\end{itemize}

\section{Threshold Calibration}

A fixed threshold does not transfer between datasets. The framework uses
a dataset adaptive strategy instead: a calibration set of known benign
documents is scored, and the threshold is set at the empirical quantile
of those scores corresponding to a chosen false positive target, for
example one percent. The dashboard exposes this target as a slider, so
the full tradeoff can be explored interactively. Figure
\ref{fig:scores} shows the intended visualization of score distributions
and the threshold.

\begin{figure}[H]
  \centering
  \figimg{scores.png}{Score distribution of benign documents with the
  scores of crafted documents marked, and the calibrated threshold line.}
  \label{fig:scores}
\end{figure}

% ======================================================================
\chapter{Implementation and Dashboard}
% ======================================================================

\section{Models and Backends}

The default embedding backend is a small sentence transformer model with
six layers. The detection study uses a small causal language model of
about 124 million parameters for perplexity scoring. Both load
automatically and are cached after first use. Selection and crafting use
pluggable LLM backends: commercial APIs when credentials are provided,
and a deterministic offline backend that simulates plausible responses so
that the entire framework remains runnable without network access or API
costs.

The gradient based crafting path runs on a GPU when available and falls
back to CPU with a warning otherwise.

\section{The Dashboard}

The dashboard exposes five views, summarized in Table
\ref{tab:dashboard}. Figure \ref{fig:dashboard} shows a screenshot of the
benchmark view.

\begin{table}[H]
  \centering
  \caption{The five dashboard views.}
  \label{tab:dashboard}
  \begin{tabular}{lp{10.8cm}}
    \toprule
    \textbf{View} & \textbf{Purpose} \\
    \midrule
    Library & View, edit, add, remove, generate, import, and export tool documents; live tool count. \\
    Benchmark & Configure backend, metric, top-$k$, and queries; inject one comparison document; run the two passes with progress; inspect metric cards, charts, and per query records. \\
    Crafting & Generate the five manual injection styles; run automated two phase crafting with the gradient free and gradient based strategies; view the resulting documents. \\
    Detection & Choose a detector, score the library and crafted documents, adjust the false positive target with live FNR, FPR, AUC, a flagged document table, and a tradeoff curve. \\
    History & Compare previously saved benchmark runs side by side and inspect any run's records. \\
    \bottomrule
  \end{tabular}
\end{table}

\begin{figure}[H]
  \centering
  \figimg{dashboard.png}{The dashboard benchmark view showing the metric
  cards, the baseline versus modified comparison chart, and the per
  query record table.}
  \label{fig:dashboard}
\end{figure}

Every expensive computation is cached. Models and API clients are loaded
once per process; embeddings and detector scores are memoized per library
state, so changing a threshold or re-running a pass recomputes only what
changed.

% ======================================================================
\chapter{Illustrative Results}
% ======================================================================

This chapter reports measurements obtained with the current synthetic
setup: the fifteen tool library, ten automatically generated weather
queries, top-$k$ equal to 5, the default sentence transformer embedder,
and the offline selection backend. These numbers illustrate the kind of
output the framework produces; they are not a large scale evaluation.

\section{Baseline and Comparison Document}

Table~\ref{tab:baseline} shows a clean baseline pass and a modified pass
in which a hand written, benign sounding alternative weather tool was
injected. The alternative document is retrieved for every query, because
it is clearly weather relevant, but it is almost never selected, because
it carries no persuasive instruction.

\begin{table}[H]
  \centering
  \caption{Baseline pass and modified pass with a hand written
  alternative weather document.}
  \label{tab:baseline}
  \begin{tabular}{lcc}
    \toprule
    \textbf{Metric} & \textbf{Baseline} & \textbf{Modified} \\
    \midrule
    Accuracy & 0.900 & n/a \\
    Hit rate at 5 & 1.000 & 1.000 \\
    Target selection rate & n/a & 0.000 \\
    Target retrieval rate & n/a & 1.000 \\
    \bottomrule
  \end{tabular}
\end{table}

\section{Injection Study}

Table~\ref{tab:injection} compares the manual injection styles with the
automatically crafted documents. The manual styles are retrieved
reasonably often but are selected rarely, because the selection backend
treats their abrupt instructions with caution. The automatically crafted
documents, whose descriptions read naturally and are tightly aligned with
the queries, achieve far higher selection rates. The gradient based
document was optimized with a short budget of five update steps; with a
larger budget the selection rate increases further.

\begin{table}[H]
  \centering
  \caption{Injection study on ten weather queries, top-5, synthetic
  library. Target selection rate and target retrieval rate per document
  style.}
  \label{tab:injection}
  \begin{tabular}{lcc}
    \toprule
    \textbf{Document style} & \textbf{Selection rate} & \textbf{Retrieval rate} \\
    \midrule
    Direct instruction & 0.100 & 0.500 \\
    Escape characters & 0.100 & 0.500 \\
    Context override & 0.100 & 0.700 \\
    Fake completion & 0.200 & 0.900 \\
    Combined & 0.200 & 0.900 \\
    Automated, gradient free & 1.000 & 1.000 \\
    Automated, gradient based & 0.500 & 0.800 \\
    \bottomrule
  \end{tabular}
\end{table}

Two observations follow. First, retrieval alone is not the whole story:
manual documents reach the candidate set but still lose at selection.
Second, effort matters: natural, query aligned text dominates abrupt
instructions, which is precisely why the two phase crafting approach
exists.

\section{Detection Study}

Table~\ref{tab:detection} reports detection results when the crafted
documents from the injection study are mixed with the fifteen benign
library documents and scored by each detector, with the threshold
calibrated at a one percent false positive target on the benign set. The
perplexity based detectors flag most of the abrupt manual documents,
whose phrasing is unusual, but consistently miss the naturally written
automated documents. The known answer detector misses essentially
everything, since the injected documents do not disturb the secret
repetition task in any way the model notices. These outcomes motivate the
central open question of the project: how to detect documents that are
adversarial without being unusual.

\begin{table}[H]
  \centering
  \caption{Detection results on the mixed set of fifteen benign and eight
  crafted documents.}
  \label{tab:detection}
  \begin{tabular}{lccc}
    \toprule
    \textbf{Detector} & \textbf{FNR} & \textbf{FPR} & \textbf{AUC} \\
    \midrule
    Perplexity & 0.125 & 0.067 & 0.97 \\
    Windowed perplexity & comparable & comparable & comparable \\
    Known answer & 1.000 & 0.000 & 0.50 \\
    \bottomrule
  \end{tabular}
\end{table}

Note that the aggregate numbers depend on the mixture: if only the
naturally written documents are scored, the perplexity detector misses
them entirely, which is the important case for a realistic adversary.

% ======================================================================
\chapter{Limitations and Future Work}
% ======================================================================

The current framework has several known limitations, each pointing to
future work.

\begin{itemize}
  \item \textbf{Scale.} The synthetic library contains fifteen tools,
  while real deployments expose hundreds to thousands. Retrieval becomes
  a real hurdle only at scale, so the retrieval half of the injection
  study should be repeated on larger generated libraries.
  \item \textbf{Backends.} Most measurements use the deterministic
  offline selection backend and small local models. Repeating the study
  with commercial LLMs and stronger retrievers would strengthen the
  findings.
  \item \textbf{Detection scope.} Only three detectors are implemented.
  Prevention style defenses, which require fine tuned model checkpoints,
  and ensemble detectors are natural additions.
  \item \textbf{Crafting budget.} The gradient based crafting path runs
  with a short update budget for interactivity. Larger budgets improve
  both document quality and success rates at the cost of compute.
\end{itemize}

% ======================================================================
\chapter{Conclusion}
% ======================================================================

This report described an evaluation framework for studying tool selection
robustness in LLM agents. The framework implements the complete two step
pipeline, supports controlled injection of crafted tool documents through
a copy on write library, measures selection and retrieval influence
separately through target rates, and evaluates three detection families
under a dataset adaptive threshold. An interactive dashboard exposes
every stage of the methodology. The illustrative results reproduce the
expected qualitative picture: hand written documents are weak, naturally
crafted documents are strong, and current detectors struggle most with
exactly the documents that matter. The framework is designed to grow
along the limitations listed above, and we welcome feedback on the
methodology, the metrics, and the priorities for the next iteration.

\end{document}
